\documentclass[tikz,border=3pt]{standalone}
\input{tikz-tcee}
\begin{document}
\hyphenpenalty=10000 \exhyphenpenalty=10000
% Cadena productividad -> CLU -> CLU reales -> margen
\begin{tikzpicture}[tcee, font=\footnotesize,
  f/.style={draw=#1, fill=#1!12, rounded corners=2pt, line width=0.7pt, minimum height=0.62cm, text width=3.6cm, align=center},
  l/.style={text width=4.1cm, align=left, font=\scriptsize},
  sig/.style={-{Stealth[length=1.6mm]}, thick, tceegris}]
  \node[f=dicproductividad] (a) at (0,0) {$\gamma_{\mathrm{PAT}}=\gamma_Y-\gamma_L$};
  \node[l, right=0.25cm of a] {\textbf{Productividad}: si el empleo cae más que el PIB, sube de forma espuria};
  \node[f=dicparo] (b) at (0,-0.95) {$\gamma_{\mathrm{CLU}}=\gamma_W-\gamma_{\mathrm{PAT}}$};
  \node[l, right=0.25cm of b] {\textbf{CLU nominales}: frente a los socios, competitividad en costes};
  \node[f=dicparo] (c) at (0,-1.9) {$\gamma_{\mathrm{CLU}/P}=\gamma_{\mathrm{CLU}}-\gamma_P$};
  \node[l, right=0.25cm of c] {\textbf{CLU reales}: participación de los salarios; $\uparrow$ presiones de costes};
  \node[f=dicpib] (d) at (0,-2.85) {$\gamma_{1+\mu}\approx\gamma_P-\gamma_{\mathrm{CLU}}$};
  \node[l, right=0.25cm of d] {\textbf{Margen}: sube exactamente cuando caen los CLU reales};
  \draw[sig] (a) -- (b); \draw[sig] (b) -- (c); \draw[sig] (c) -- (d);
\end{tikzpicture}
\end{document}
